\section{Crystal form factors}\label{app:E.7}

An electron-recoil rate of silicon at independent-particle level is a sum of $|\rho_{cv}(\vect{k},\vect{q})|^2$ over valence states $v\vect{k}$, conduction states $c\vect{k}'$ and momentum transfers $\vect{q}$, each term weighted by an energy delta function. With four valence bands, 52 conduction bands, a $10\times10\times10$ grid of crystal momenta and every reciprocal-lattice vector up to $q = \SI{30}{keV}$, the sum has about $5\times10^{11}$ terms (\cref{we:E-siff}). Each term is an overlap of two sets of plane-wave coefficients. The work is in restoring the core region that the pseudopotential removed, adding the core shells, and knowing when the sum has converged; the bands come from the calculation of \cref{app:E.1}, and the screening that turns the bare vertex into a measured loss function from \cref{app:E.6}.

\subsection{The density vertex}

The object is the density vertex of \cref{eq:A.5-rho},
\begin{equation*}
  \rho_{mn}(\vect{k},\vect{q}) = \langle\psi_{m,\vect{k}+\vect{q}}|\ee^{\ii\vect{q}\cdot\vect{r}}|\psi_{n\vect{k}}\rangle ,
\end{equation*}
with $\vect{q}$ never reduced to the first zone (\cref{app:A.3,app:A.4}). From it the independent-particle structure factor follows as a sum over pairs (\cref{sec:10.1}); the crystal form factor $f_{\mathrm{crys}}$ of the literature is the same object in another normalization, \cref{eq:A.8-fcrys}; and the vertices of the spin and current couplings of \cref{sec:10.4} are built the same way, with $\ee^{\ii\vect{q}\cdot\vect{r}}$ multiplied by a Pauli matrix or by the velocity.

\subsection{Place on the hierarchy}

On the truncation label $\mathfrak{T}$ of \cref{sec:8.5}, a form-factor calculation sits at the base point $\mathfrak{T}_0$. The nuclei are clamped at their equilibrium positions, $n_B = 0$; there is no electron--phonon insertion, $n_g = 0$, and no anharmonic or thermal correction, $n_A = 0$ and zeroth thermal order, $n_T = 0$; the many-body approximation is that of independent particles, $\mathcal{X} = \mathrm{IP}$. The plane wave is kept whole, so the result is exact to all orders in the multipole parameter $\xi_D$, and the dipole form $\vect{q}\cdot\vect{v}_{mn}/(\varepsilon_{m\vect{k}} - \varepsilon_{n\vect{k}})$ of \cref{eq:A.5-optical} is only a check at small $q$. The Kohn--Sham approximation and the pseudopotential sit inside $\mathcal{X} = \mathrm{IP}$ and are not part of the label. The states are Kohn--Sham states rather than quasiparticles (\cref{sec:3.4}), with a scissor shift $\Delta_{\mathrm{scis}}$ on the conduction energies or with $GW$ energies (\cref{app:E.3}); and they are pseudo-wavefunctions unless the region inside the core radius is reconstructed. What the base point leaves out is screening and local fields ($\mathcal{X} = \mathrm{RPA}$ and beyond, \cref{app:E.6}), excitons ($\mathcal{X} = \mathrm{BSE}$), and the lattice ($n_g \ge 1$, \cref{app:E.5,app:E.8}). Screening divides the squared potential by $|\varepsilon(\vect{q},\omega)|^2$, which is $\varepsilon_\infty^2$ as $q \to 0$ and approaches one only at large $q$ (\cref{sec:4.4}), so a form factor at $\mathfrak{T}_0$ is an input to that correction, not a substitute for it.

\subsection{Plane-wave evaluation}

A plane-wave code returns, for each band $n$ and each $\vect{k}$ of its grid, the coefficients $c_{n\vect{k}}(\vect{G})$ of
\begin{equation}
  u_{n\vect{k}}(\vect{r}) = \frac{1}{\sqrt\Omega}\sum_{|\vect{k}+\vect{G}|^2\le 2m_eE_{\mathrm{cut}}} c_{n\vect{k}}(\vect{G})\,\ee^{\ii\vect{G}\cdot\vect{r}}, \qquad \sum_{\vect{G}}|c_{n\vect{k}}(\vect{G})|^2 = 1 ,
  \label{eq:E.7-pw}
\end{equation}
which is the normalization of \cref{eq:A.5-bloch} and the code convention of \cref{app:A.3} (coefficients of $\ee^{+\ii\vect{G}\cdot\vect{r}}$, the opposite of the book's lattice Fourier components). The cutoff $E_{\mathrm{cut}}$ is the one of \cref{sec:3.5}. With $\vect{k}+\vect{q}$ folded back into the zone, $\vect{k}+\vect{q} = \vect{k}' + \vect{G}_{\mathrm{U}}$ with $\vect{k}'$ a point of the grid and $\vect{G}_{\mathrm{U}}$ a reciprocal-lattice vector, and by \cref{eq:A.5-periodic} the vertex is $\int_\Omega u^*_{m\vect{k}'}\ee^{\ii\vect{G}_{\mathrm{U}}\cdot\vect{r}}u_{n\vect{k}}$, and inserting \cref{eq:E.7-pw} the cell integral of $\ee^{\ii(\vect{G}_{\mathrm{U}} + \vect{G} - \vect{G}')\cdot\vect{r}}$ is $\Omega\,\delta_{\vect{G}',\vect{G}+\vect{G}_{\mathrm{U}}}$:
\begin{equation}
  \rho_{mn}(\vect{k},\vect{q}) = \sum_{\vect{G}} c^*_{m\vect{k}'}(\vect{G} + \vect{G}_{\mathrm{U}})\,c_{n\vect{k}}(\vect{G}), \qquad \vect{q} = \vect{k}' - \vect{k} + \vect{G}_{\mathrm{U}} .
  \label{eq:E.7-rhopw}
\end{equation}
The sum runs over the vectors present in both basis sets. Coefficients are matched by their Miller indices, the integer coordinates of $\vect{G}$ in the basis $\vect{b}_i$, since the two states are stored on different lists of $\vect{G}$. A term with $\vect{G}+\vect{G}_{\mathrm{U}}$ outside the basis of $\vect{k}'$ is zero in the representation, so \cref{eq:E.7-rhopw} is exact for the functions the code represents; its error is the basis truncation itself, which we discuss below. This is the expression of Ref.~\cite[Eq.~(3.16)]{essig2016direct}, in the notation of \cref{app:A.5}.

Each value of \cref{eq:E.7-rhopw} costs one dot product of length $N_{\mathrm{PW}}$, the number of plane waves per state. All the $\vect{G}_{\mathrm{U}}$ for one pair $(m\vect{k}', n\vect{k})$ come at once from a fast Fourier transform: the product $u^*_{m\vect{k}'}(\vect{r})\,u_{n\vect{k}}(\vect{r})$ is formed on the real-space grid and transformed, and the component at $\vect{G}_{\mathrm{U}}$ is $\rho_{mn}$ for $\vect{q} = \vect{k}' - \vect{k} + \vect{G}_{\mathrm{U}}$. The cost per pair drops from $N_{\mathrm{PW}}N_{\vect{G}}$ to of order $N_{\mathrm{FFT}}\log N_{\mathrm{FFT}}$, where $N_{\vect{G}}$ is the number of $\vect{G}_{\mathrm{U}}$ needed and $N_{\mathrm{FFT}}$ the number of grid points~\cite[below Eq.~(3.15)]{griffin2021extended}. The real-space grid has to resolve the product, whose components reach twice the wavefunction cutoff. For a cubic cell it covers a cube of side $4\sqrt{2m_eE_{\mathrm{cut}}}$ in reciprocal space, which holds $8\Omega(2m_eE_{\mathrm{cut}})^{3/2}/\pi^3$ points, while the sphere of \cref{eq:E.7-pw} holds $N_{\mathrm{PW}} = \Omega(2m_eE_{\mathrm{cut}})^{3/2}/(6\pi^2)$, so $N_{\mathrm{FFT}} \approx (48/\pi)N_{\mathrm{PW}}$. The fcc grid runs along the primitive vectors, whose lengths multiply to $\sqrt2\,\Omega$, and holds $\sqrt2$ times more, $N_{\mathrm{FFT}} \approx (48\sqrt2/\pi)N_{\mathrm{PW}} \approx 22\,N_{\mathrm{PW}}$: about $36^3$ points at \SI{60}{Ry}, as in \cref{app:E.1}.

\paragraph{Which $\vect{q}$ a grid can reach.} \Cref{eq:E.7-rhopw} needs the state at $\vect{k}' = \vect{k}+\vect{q} - \vect{G}_{\mathrm{U}}$. If both $\vect{k}$ and $\vect{k}'$ are taken from one uniform grid of $n_k = n_1n_2n_3$ points (\cref{app:C.2}), the first-zone part of $\vect{q}$ is a difference of grid vectors, and $\vect{q}$ is resolved in steps of $|\vect{b}_i|/n_i$. A form factor tabulated on a finer or arbitrary set of $\vect{q}$ needs a second, shifted set of states: one non-self-consistent calculation at the points $\vect{k}+\vect{q}$, as density-functional perturbation theory does for phonon wavevectors (\cref{app:E.4}). The states at $\vect{k}'$ must be computed, or unfolded from the irreducible wedge by a symmetry operation, with the phases that the operation carries. In the diamond structure some operations include a fractional translation $\vect{t}_{\mathrm{f}}$, and each coefficient of a rotated Bloch state picks up a phase $\ee^{-\ii\vect{k}_{\mathrm{rot}}\cdot\vect{t}_{\mathrm{f}}}$, with $\vect{k}_{\mathrm{rot}}$ the rotated wavevector $\vect{k}+\vect{G}$. Computing the full grid without symmetry avoids the bookkeeping at a cost of the order of the symmetry factor, 48 for silicon.

\paragraph{Gauge and degeneracy.} Each Bloch state carries an arbitrary phase, and $\rho_{mn}$ picks up the phase difference of its two states, while $|\rho_{mn}|^2$ does not. A rate that sums $|\rho_{mn}|^2$ is therefore gauge independent, provided that the sum over $m$ and $n$ covers complete degenerate sets: a degenerate subspace may be rotated by any unitary matrix, and only the trace over it is fixed. At the valence maximum of silicon the threefold degeneracy at $\mathrm{\Gamma}$ makes this concrete; a sum restricted to one of the three states there is meaningless. Amplitudes that interfere, such as the second-order amplitudes of phonon-assisted processes in which $\rho_{mn}$ multiplies an electron--phonon vertex, need $\rho$ and $g$ in one gauge, computed from the same stored states (\cref{app:A.5,app:E.5}).

\subsection{Velocity matrix elements and nonlocal potentials}

The momentum matrix element follows from \cref{eq:E.7-pw} in the same way,
\begin{equation}
  \vect{p}_{mn}(\vect{k}) = \sum_{\vect{G}} c^*_{m\vect{k}}(\vect{G})\,(\vect{k} + \vect{G})\,c_{n\vect{k}}(\vect{G}) ,
  \label{eq:E.7-ppw}
\end{equation}
but the velocity is not $\vect{p}_{mn}/m_e$ in a pseudopotential calculation. The velocity operator is $\vect{v} = \ii[H,\vect{r}]$. For $H = \vect{p}^2/2m_e + V_{\mathrm{loc}}(\vect{r}) + V_{\mathrm{NL}}$, with $V_{\mathrm{NL}}$ the nonlocal part of the pseudopotential (\cref{app:A.5}), the local potential commutes with $\vect{r}$ and
\begin{equation}
  \vect{v} = \frac{\vect{p}}{m_e} + \ii[V_{\mathrm{NL}},\vect{r}] .
  \label{eq:E.7-v}
\end{equation}
The commutator is evaluated in the $\vect{k}$ representation. Define $V_{\mathrm{NL}}(\vect{k}) \equiv \ee^{-\ii\vect{k}\cdot\vect{r}}V_{\mathrm{NL}}\ee^{\ii\vect{k}\cdot\vect{r}}$, the operator that acts on $u_{n\vect{k}}$. Differentiating, $\nabla_{\vect{k}}V_{\mathrm{NL}}(\vect{k}) = \ee^{-\ii\vect{k}\cdot\vect{r}}\,\ii[V_{\mathrm{NL}},\vect{r}]\,\ee^{\ii\vect{k}\cdot\vect{r}}$, so that
\begin{equation}
  \vect{v}_{mn}(\vect{k}) = \frac{\vect{p}_{mn}(\vect{k})}{m_e} + \langle u_{m\vect{k}}|\nabla_{\vect{k}}V_{\mathrm{NL}}(\vect{k})|u_{n\vect{k}}\rangle .
  \label{eq:E.7-vNL}
\end{equation}
For the separable form \cref{eq:E.2-KB} that plane-wave codes use, with projectors $\tilde p_{\kappa i}$ centered on the atoms and coefficients $D_{\kappa,ij}$, the gradient acts only on the Fourier transforms of the projectors, $\tilde p_{\kappa i}(\vect{k}+\vect{G})$, and costs about as much as applying $V_{\mathrm{NL}}$ once. The correction is not small. In diamond, the absorption spectrum computed with $\vect{p}/m_e$ instead of $\vect{v}$ differs by 15 to 20 percent~\cite[Fig.~1]{ismailbeigi2001coupling}, and Ref.~\cite[Eq.~(20)]{ismailbeigi2001coupling} shows that a long-wavelength field couples to $\vect{v}$, not to $\vect{p}/m_e$.

For the dark-matter vertices the correction has opposite consequences for the density and for the current. The density vertex needs none. The operator $\ee^{\ii\vect{q}\cdot\vect{r}}$ is local, and \cref{eq:E.7-rhopw} evaluated at small $q$ reproduces $\vect{q}\cdot\vect{v}_{mn}/(\varepsilon_{m\vect{k}} - \varepsilon_{n\vect{k}})$ with the full velocity of \cref{eq:E.7-vNL}, since the derivation of \cref{eq:A.5-optical} used only $\ii[H,\vect{r}] = \vect{v}$. The comparison of the two sides at $q \sim 10^{-3}\,|\vect{b}_i|$ is a test that catches a missing nonlocal term in a velocity routine. A coupling to the electron current, on the other hand, is ambiguous with a nonlocal Hamiltonian. The paramagnetic current $j_e(\vect{q})$ of \cref{sec:4.3} is built from $\vect{p}/m_e$; gauge invariance requires an additional nonlocal term, which Ref.~\cite[Eqs.~(18) and~(19)]{ismailbeigi2001coupling} derives as a phase $\exp(\ii\int\vect{A}\cdot\dd\vect{x})$ along the straight line between the two arguments of $V_{\mathrm{NL}}(\vect{r},\vect{r}')$. Its long-wavelength limit is the commutator of \cref{eq:E.7-v}. At finite $q$ the term depends on $\vect{q}$ through the extent of the projectors, of order $q\,r_c$, and codes that compute current responses at finite $q$ either include it or restrict themselves to $q\,r_c \ll 1$.

\paragraph{The scissor shift.} The scissor operator $\Delta_{\mathrm{scis}}\sum_c|\psi_c\rangle\langle\psi_c|$ is nonlocal too. It rescales each interband velocity by $(\varepsilon_{c\vect{k}} - \varepsilon_{v\vect{k}} + \Delta_{\mathrm{scis}})/(\varepsilon_{c\vect{k}} - \varepsilon_{v\vect{k}})$, the rigid-shift case of \cref{eq:E.3-vrescale}, and leaves the density vertex and its optical limit unchanged. An optical spectrum built from unscaled velocities acquires the spurious factor of \cref{app:E.6}; first-principles absorption calculations therefore carry the nonlocality of the scissor into the oscillator strengths~\cite[Sect.~VII]{zacharias2016one}, where the inverse factor appears because it multiplies their oscillator-strength form rather than the velocity. For silicon in the local-density approximation, with $\Delta_{\mathrm{scis}} = \SI{0.74}{eV}$~\cite[Sect.~VII]{zacharias2016one} and a direct gap at $\mathrm{\Gamma}$ of \SI{2.57}{eV}~\cite[Table~V]{ponce2015temperature}, the factor is 1.29 at $\mathrm{\Gamma}$, and $|\vect{v}_{cv}|^2$ changes by 66 percent.

\subsection{All-electron reconstruction}

A pseudo-wavefunction agrees with the all-electron one outside a sphere of radius $r_c$ around each nucleus and is smooth and nodeless inside it (\cref{sec:3.5}). The two differ only inside the sphere, so their difference has Fourier components at momenta of order $\pi/r_c$ and above, and the pseudo vertex is reliable only for $q\,r_c \ll 1$. The PAW map \cref{eq:E.2-PAW} restores the difference, in one of two ways. The first route expands the all-electron state itself in plane waves up to a second, larger cutoff and then applies \cref{eq:E.7-rhopw} unchanged. Ref.~\cite[Sect.~2.1.1]{griffin2021extended} does this for silicon and germanium with a second cutoff of \SI{2}{keV}, which by $q_{\mathrm{max}} = (2m_eE)^{1/2}$ reaches $q \approx \SI{45}{keV}$. The second route works with the matrix elements, \cref{eq:E.2-rhoPAW}. With $\vect{r} = \vect{\tau}_\kappa + \vect{x}$ in each sphere, the Rayleigh expansion \cref{eq:C.3-planewave} of $\ee^{\ii\vect{q}\cdot\vect{x}}$ reduces each element of the one-center matrix \cref{eq:E.2-onecentre} to Gaunt coefficients times the radial integrals
\begin{equation}
  \int_0^{r_c}\dd x\,x^2\,j_l(qx)\,\bigl[\phi^{\mathrm{AE}}_{\kappa i}(x)\phi^{\mathrm{AE}}_{\kappa j}(x) - \phi^{\mathrm{PS}}_{\kappa i}(x)\phi^{\mathrm{PS}}_{\kappa j}(x)\bigr] ,
  \label{eq:E.7-radial}
\end{equation}
with $\phi(x)$ the radial parts of the partial waves of \cref{eq:E.2-PAW} and $j_l$ the spherical Bessel functions of \cref{app:C.3}. The integrals live on the radial grid of the atomic calculation, which resolves the nodes of the all-electron partial waves, so the correction has no plane-wave cutoff at all; it is limited by the completeness of the partial waves, usually two per angular momentum. They are computed once per species and per $q$, and the per-state work is in the projections $\langle\tilde p_{\kappa j}|\tilde\psi_{n\vect{k}}\rangle$, which the code already stores.

How much the reconstruction matters depends on which $q$ the rate weights. For a heavy mediator, $\FDM = 1$, and the factor $q$ of the measure in \cref{eq:A.8-rateiso} gives large momentum transfers more weight than any lighter mediator does, and the reconstruction raises the silicon rate even above low thresholds (\cref{app:E.2}). For a light mediator, $\FDM^2 \propto q^{-4}$ weights the smallest momentum transfers, and the reconstruction changes the rate little except at high thresholds, where only large $q$ is kinematically allowed. Norm-conserving pseudopotentials can be reconstructed in the same way when the generator writes out its all-electron and pseudo partial waves.

\subsection{Core and semicore states}

Below the valence bands of silicon lie the atomic shells $2p$, $2s$ and $1s$, with measured binding energies relative to the valence maximum of \SI{99.4}{eV} ($2p_{3/2}$; $2p_{1/2}$ at \SI{99.8}{eV}), \SI{150}{eV} and \SI{1839}{eV}; in germanium the $3d$ shell sits only \SI{29}{eV} below it~\cite[Table~1-1]{williams2009electron}. Kinematics decides when they matter. A transition from a shell at depth $E_{\mathrm{core}}$ needs $\omega \ge \Eg + E_{\mathrm{core}}$, and since $\vmin = \omega/q + q/2m_\chi \le v_{\mathrm{max}}$ it also needs $q \ge \omega/v_{\mathrm{max}}$ and $m_\chi \ge 2\omega/v_{\mathrm{max}}^2$ (\cref{sec:1.3}). With $v_{\mathrm{max}} = 2.650\times10^{-3}$, the silicon $2p$ shell opens at $m_\chi = \SI{29}{MeV}$ and needs $q \ge \SI{38}{keV}$, and the $2s$ shell at \SI{43}{MeV} with $q \ge \SI{57}{keV}$; the germanium $3d$ shell opens at \SI{8.5}{MeV} and needs $q \ge \SI{11}{keV}$. Silicon $2p$ transitions therefore probe the wavefunctions at momenta beyond the reach of a typical plane-wave basis (\cref{sec:3.5}), and a pseudopotential calculation does not contain the shell at all.

Such shells are either kept in the pseudopotential or taken from the atom. A \emph{semicore} shell, such as germanium $3d$, can be kept in the valence of the pseudopotential, at the price of a higher cutoff, and is then computed with the other bands and reconstructed as above. A deeper shell is taken from the isolated atom. Its Bloch states are the lattice sums
\begin{equation}
  \psi_{\kappa nlm,\vect{k}}(\vect{r}) = \frac{1}{\sqrt N}\sum_{\vect{R}}\ee^{\ii\vect{k}\cdot(\vect{R}+\vect{\tau}_\kappa)}\,\psi^{\mathrm{at}}_{\kappa nlm}(\vect{r} - \vect{R} - \vect{\tau}_\kappa) ,
  \label{eq:E.7-core}
\end{equation}
of the atomic orbitals $\psi^{\mathrm{at}}_{\kappa nlm}$~\cite[Eq.~(2.8)]{griffin2021extended}, labeled by the atomic quantum numbers instead of a band index; the orbitals barely overlap, so the band is flat and its energy is the atomic level. The atomic approximation can be tested where both treatments exist. Ref.~\cite[Sect.~2.2]{griffin2021extended} computes the silicon $2p$ and germanium $3d$ shells both ways and finds that the atomic orbitals reproduce the reconstructed crystal wavefunctions up to the largest momentum the crystal calculation resolves, about \SI{45}{keV}. The vertex between a core state and a conduction state is then $\int\dd^3r\,\psi^*_{c,\vect{k}+\vect{q}}\ee^{\ii\vect{q}\cdot\vect{r}}\psi_{\kappa nlm,\vect{k}}$, with the conduction state from the crystal calculation; above the highest computed band the final states are approximated by plane waves, in the spirit of \cref{eq:2.2-fion-pw}, with or without a Coulomb factor~\cite[Sects.~2.3 and~3.4]{griffin2021extended}. Relativistic corrections to the core orbitals are of relative size $(Z\alpha)^2$, about \SI{1}{\percent} for the silicon $1s$ shell and \SI{5}{\percent} for germanium $1s$. That is the size of each term; the net shift of the $1s$ level is about a quarter of it, \SI{0.26}{\percent} for silicon and \SI{1.46}{\percent} for germanium (\cref{we:3-Zalpha}). The approximation leaves out the relaxation of the other electrons around the core hole, which shifts and reshapes the edge; that final-state effect lies beyond $\mathcal{X} = \mathrm{IP}$.

\subsection{Atomic ionization form factors}

For a target made of isolated atoms, such as a noble liquid, the rate of \cref{sec:2.2} needs the ionization form factor $|f_{\mathrm{ion}}(p_{\mathrm{ej}},q)|^2$ of \cref{eq:2.2-fion} for each occupied shell. The initial orbitals are Hartree--Fock orbitals (\cref{sec:3.4}), almost always the Roothaan--Hartree--Fock solutions of Bunge and collaborators~\cite{bunge1993roothaan}, which expand each radial function in Slater-type orbitals,
\begin{equation}
  R^{\mathrm{at}}_{nl}(r) = a_0^{-3/2}\sum_j C_{jln}\,\frac{(2Z_{jl})^{\tilde n_{jl}+1/2}}{\sqrt{(2\tilde n_{jl})!}}\Bigl(\frac{r}{a_0}\Bigr)^{\tilde n_{jl}-1}\ee^{-Z_{jl}r/a_0} ,
  \label{eq:E.7-sto}
\end{equation}
with tabulated coefficients $C_{jln}$, exponents $Z_{jl}$ and integers $\tilde n_{jl}$~\cite[Eq.~(B33)]{catena2020atomic}. The final state is a continuum state of energy $E_{\mathrm{ej}} = p_{\mathrm{ej}}^2/2m_e$, expanded in partial waves $l'$. In the simplest and most common choice it is the exact Coulomb wave of a hydrogenic potential $-Z_{\mathrm{eff}}\alpha/r$, with $Z_{\mathrm{eff}} = n\,(E_B/\SI{13.6}{eV})^{1/2}$ fixed by the binding energy $E_B$ of the shell~\cite[Eqs.~(B34) and~(B35)]{catena2020atomic}; a better choice solves the radial equation in the self-consistent potential of the ion. The radial integrals of $j_L(qr)$ between the two radial functions are then evaluated for every $(l, l', L)$ allowed by the Gaunt coefficients.

The convergence parameters are the partial-wave cutoff $l'_{\mathrm{max}}$, the grids in $q$ and $E_{\mathrm{ej}}$, and the accuracy of the radial integrals. Ref.~\cite[footnote~11]{catena2020atomic} sums to $l'_{\mathrm{max}} = 7$; the number needed grows with $q$, since a plane wave of momentum $q$ reaches angular momenta up to about $q$ times the orbital radius, and has to be checked at the largest $q$ of the table. The $q$ grid must extend to the kinematic maximum $2m_\chi v_{\mathrm{max}}$ of the heaviest mass considered, \SI{5.3}{MeV} at \SI{1}{GeV}, and is best logarithmic. The radial integrands oscillate rapidly at large $q$, and an integration split into intervals can misread an accidental cancellation over one interval as convergence; Ref.~\cite[Appendix~B\,5]{catena2020atomic} therefore requires two consecutive intervals to agree and switches to a series representation at large $q$. The cost is negligible next to any crystal calculation: a few shells, a few partial waves, and one-dimensional integrals.

The approximation behind the table is that of frozen orbitals. The binding energies are single-particle (Koopmans) energies, the other electrons do not relax when one is removed, and the final state does not feel the polarization of the ion. It is also nonrelativistic, which matters for the inner shells of xenon, where $(Z\alpha)^2 \approx 0.16$. A subtler problem is orthogonality. A continuum state of the model potential $-Z_{\mathrm{eff}}\alpha/r$ is not an eigenstate of the Hartree--Fock Hamiltonian and is not orthogonal to the bound orbital, so that $f_{\mathrm{ion}}$ does not vanish as $q \to 0$, as it must for $\langle f|i\rangle = 0$. Orthogonalizing the continuum state to the occupied orbitals removes the spurious small-$q$ weight, which a light mediator would otherwise amplify. For silicon these atomic tables play a second role: the deep shells of the crystal in \cref{eq:E.7-core} are built from the same orbitals.

\subsection{Convergence parameters and checks}

A crystal form factor is controlled by the cutoff, the number of conduction bands, the $\vect{k}$ grid and the treatment of the energy delta function. The cutoff $E_{\mathrm{cut}}$ sets the largest momentum the vertex contains, $(2m_eE_{\mathrm{cut}})^{1/2}$, which is \SI{31}{keV} at \SI{70}{Ry}; Ref.~\cite[Appendix~C]{essig2016direct} varies it from 30 to \SI{110}{Ry} and finds the silicon rate at $m_\chi = \SI{1}{GeV}$ changed by less than one percent above \SI{70}{Ry}. The number of conduction bands $N_c$ sets the largest energy deposition the spectrum contains; with 52 conduction bands the lowest band left out lies about \SI{57}{eV} above the gap, and the computed spectrum ends there~\cite[Sect.~4.2]{essig2016direct}. A spectrum above the highest band is zero, not small, and a rate integrated over $\omega$ is missing everything beyond it. The $\vect{k}$ grid sets both the resolution in $q$ and the sampling of the onset near $E_{\mathrm{min}}(q)$; Ref.~\cite[Appendix~C]{essig2016direct} finds the grid matters only at low masses and low thresholds, where few transitions are kinematically allowed. It tests the grid by shaking every grid point randomly by up to half the grid spacing and comparing ten such grids. Finally, the delta function of energy needs a broadening or a tetrahedron integration, with the width tied to the grid spacing as described in \cref{app:C.2}, or a binning: Ref.~\cite[Sect.~4.2]{essig2016direct} uses bins of \SI{0.1}{eV} in $\omega$ and $0.02\,\alpha m_e = \SI{75}{eV}$ in $q$.

Most errors are caught by the small-$q$ limit and by a sum rule. The small-$q$ check is the normalization $\rho_{nn}(\vect{k},0) = 1$ with $\rho_{mn}(\vect{k},\vect{q})$ for $m \neq n$ vanishing linearly in $q$, with the slope given by the velocity, as discussed above. The sum rule holds for each occupied state and measures the effect of truncating the bands. Since $\rho_{mn}(\vect{k},\vect{q}) = \langle u_{m,\vect{k}+\vect{q}}|u_{n\vect{k}}\rangle$ (\cref{eq:A.5-rho}), and the $u_{m,\vect{k}+\vect{q}}$ form a complete set of eigenstates of $H(\vect{k}+\vect{q}) = \ee^{-\ii(\vect{k}+\vect{q})\cdot\vect{r}}H\ee^{\ii(\vect{k}+\vect{q})\cdot\vect{r}}$,
\begin{equation}
  \sum_m(\varepsilon_{m,\vect{k}+\vect{q}} - \varepsilon_{n\vect{k}})\,|\rho_{mn}(\vect{k},\vect{q})|^2 = \langle u_{n\vect{k}}|H(\vect{k}+\vect{q}) - H(\vect{k})|u_{n\vect{k}}\rangle
  = \vect{q}\cdot\vect{v}^{\mathrm{g}}_n(\vect{k}) + \frac{q^2}{2m_e} + \delta_{\mathrm{NL}} ,
  \label{eq:E.7-fsum}
\end{equation}
where the sum over $m$ runs over all bands and $\delta_{\mathrm{NL}}$ is the part of $\langle u_{n\vect{k}}|V_{\mathrm{NL}}(\vect{k}+\vect{q}) - V_{\mathrm{NL}}(\vect{k})|u_{n\vect{k}}\rangle$ beyond first order in $q$. The kinetic part of $H(\vect{k}) = (\vect{p}+\vect{k})^2/2m_e + V_{\mathrm{loc}} + V_{\mathrm{NL}}(\vect{k})$ gives $\vect{q}\cdot(\vect{p}+\vect{k})/m_e + q^2/2m_e$, and the first-order part of the nonlocal term combines with it into the group velocity, by the Hellmann--Feynman relation $\nabla_{\vect{k}}\varepsilon_{n\vect{k}} = \langle u_{n\vect{k}}|\nabla_{\vect{k}}H(\vect{k})|u_{n\vect{k}}\rangle$. Summed over a filled band the linear term cancels. Transitions between two occupied states cancel in pairs, since time reversal gives $|\rho_{nm}(-\vect{k}-\vect{q},\vect{q})| = |\rho_{mn}(\vect{k},\vect{q})|$ with the opposite energy difference. What remains is the valence $f$-sum rule of a pseudo-Hamiltonian, the band-summed form of \cref{eq:E.2-dc}: $q^2/2m_e$ per valence electron, plus the nonlocal term, which is absent for the all-electron Hamiltonian of \cref{eq:A.7-fsum}. The right side of \cref{eq:E.7-fsum} is a ground-state expectation value and costs nothing; the left side with $N_c$ conduction bands falls short of it by the weight of the bands left out. With a nonlocal pseudopotential, a test that compares the truncated sum with $q^2/2m_e$ alone confuses the two effects, and the nonlocal term must be included before the shortfall is read as band truncation. The shortfall falls slowly with $N_c$, because at momentum transfer $q$ much of the missing weight sits near the free-electron energy $q^2/2m_e$.

\subsection{Cost and pitfalls}

The count of vertices sets the cost. With both $\vect{k}$ and $\vect{k}'$ on a grid of $n_k$ points, $N_v$ valence and $N_c$ conduction bands, and $N_{\vect{G}}$ reciprocal-lattice vectors within the largest $q$, there are $n_k^2N_vN_cN_{\vect{G}}$ vertices, computed with $n_k^2N_vN_c$ fast Fourier transforms of $N_{\mathrm{FFT}}$ points each. Directional averages can be taken over the irreducible wedge of $\vect{k}$, with the full set of $\vect{k}'$. The vertices themselves are far too many to store, and they are accumulated into bins of $(q,\omega)$, or of $(\vect{q},\omega)$ for a directional response, as they are computed; a tabulated crystal form factor is that binned sum. The all-electron correction adds, per state, the projections onto the atomic projectors, and per $q$ the radial integrals of \cref{eq:E.7-radial}.

Three errors are not caught by the checks above. A $\vect{q}$ reduced to the first zone assigns an Umklapp vertex to the wrong momentum transfer (\cref{app:A.3}). A spectrum cut by the highest band, or a $q$ range cut by $E_{\mathrm{cut}}$, looks converged in its own window and silently loses everything outside it. And normalizations differ between codes and papers: plane-wave codes often fold the spin factor 2 into the $\vect{k}$ weights~\cite[Eq.~(4.2)]{essig2016direct}, and the conversions between $f_{\mathrm{crys}}$ and $\bar S$ are in \cref{eq:A.8-fcrys}. For germanium, which semilocal functionals leave without a gap~\cite[Sect.~2.1]{griffin2021extended}, the band order must be fixed before any rigid shift (\cref{app:E.1}).

\begin{workedexample}{The silicon form factor}\label{we:E-siff}
We set up a valence-to-conduction calculation for silicon at $\mathfrak{T}_0$ and count its cost. The cell volume is $\Omega = a^3/4 = \SI{40.05}{\angstrom^3}$ at $a = \SI{5.431}{\angstrom}$. A cutoff of \SI{70}{Ry}, \SI{952}{eV}, reaches $(2m_eE_{\mathrm{cut}})^{1/2} = \SI{31.2}{keV} = \SI{15.8}{\angstrom^{-1}}$ and holds $N_{\mathrm{PW}} = \Omega(2m_eE_{\mathrm{cut}})^{3/2}/6\pi^2 \approx 2700$ plane waves per state. The non-self-consistent run computes 4 valence and 52 conduction bands on a $10\times10\times10$ grid without symmetry, $n_k = 1000$; the coefficients take $1000\times56\times2700$ complex numbers, \SI{2.4}{GB} in double precision, which fits in memory. The Kohn--Sham gap is about \SI{0.6}{eV} in the representative PBE calculation (\cref{we:E-si-bands}), and a scissor $\Delta_{\mathrm{scis}} = \Eg - E_{\mathrm{g}}^{\mathrm{KS}}$ moves the conduction bands up to the measured low-temperature gap, \SI{1.17}{eV} (\cref{sec:1.3}).

The reciprocal-lattice vectors within $q_{\mathrm{max}} = \SI{30}{keV}$ number $\frac{4\pi}{3}q_{\mathrm{max}}^3\Omega/(2\pi)^3 \approx 2400$, since the zone of silicon has volume $(2\pi)^3/\Omega = \SI{47.6}{keV^3}$. The vertex count is $n_k^2N_vN_cN_{\vect{G}} = 10^6\times4\times52\times2400 \approx 5\times10^{11}$, the number quoted at the start of this section. They come from $n_k^2N_vN_c = 2.1\times10^8$ transforms on grids of $N_{\mathrm{FFT}} \approx 22N_{\mathrm{PW}} \approx 6\times10^4$ points. At about $5N_{\mathrm{FFT}}\log_2N_{\mathrm{FFT}} \approx 5\times10^6$ floating-point operations per transform, the total is $1\times10^{15}$ operations, about three hours on a node that sustains $10^{11}$ operations per second, and less with the irreducible wedge. The calculation is cheap; the binned output, at the bin widths quoted above, has about $500\times400$ entries up to \SI{30}{keV} (the published table of Ref.~\cite{essig2016direct} has $500\times900$ and reaches \SI{67}{keV}).

The calculation is then tested at small $q$, against the sum rule and against the cutoff. At $q = 10^{-3}\,|\vect{b}_i|$ the ratio $|\rho_{cv}|/q$ must equal $|\hat{\vect{q}}\cdot\vect{v}_{cv}|/(\varepsilon_c - \varepsilon_v)$ with the velocity of \cref{eq:E.7-vNL}. For each valence state the truncated left side of \cref{eq:E.7-fsum}, with its nonlocal term on the right, measures the band truncation. And the rate at $m_\chi = \SI{1}{GeV}$ for a heavy mediator, recomputed at a higher cutoff and with all-electron reconstruction, fixes where the pseudo vertex stops being adequate. Kinematics tells us where to look: at \SI{10}{MeV}, with $\omega$ just above the gap, momentum transfers between \SI{0.45}{keV} and \SI{53}{keV} are allowed, but most of the valence weight lies at a few keV, and the cutoff test decides how much of the rate comes from beyond \SI{30}{keV}; the $2p$ shell, which would need $q \ge \SI{38}{keV}$, is still closed below \SI{29}{MeV}.
\end{workedexample}
